\documentclass[10pt,a4paper]{amsart}
\usepackage[margin=19mm]{geometry}
\usepackage{amsmath,amssymb,mathtools}
\usepackage[scr=rsfs,cal=euler]{mathalfa}
\usepackage{xfrac}
\usepackage[unicode,hidelinks,bookmarks=false]{hyperref}
\newcommand{\ie}{i.e.\ }
\newcommand{\cf}{cf.\ }
\newcommand{\resp}{respectively\ }
\newcommand{\Subsection}{\subsection}
\providecommand{\nobreakdash}{\nobreak\hbox{-}}
\setlength{\emergencystretch}{2em}
\multlinegap0pt
\usepackage{amssymb}
\usepackage{mathtools}
\usepackage[shortlabels]{enumitem}
\newtheorem{lemma}{Lemma}
\newtheorem{theorem}{Theorem}
\newcommand{\D}{\mathbb D}
\newcommand{\dd}{\mathrm d}
\newcommand{\ii}{\mathrm i}
\title{The Nitsche--Hopf Conjecture for Minimal Graphs}
\author{David Kalaj, Jian-Feng Zhu}
\date{Source: arXiv:2605.10140v2 (CC BY 4.0)}
\begin{document}
\maketitle

\noindent\emph{Source and licence.} The Nitsche--Hopf Conjecture for Minimal Graphs. Version arXiv:2605.10140v2 (CC BY 4.0), distributed by arXiv under the Creative Commons Attribution 4.0 International licence, http://creativecommons.org/licenses/by/4.0/, as recorded in the arXiv licence metadata for this exact version and shown on the abstract page https://arxiv.org/abs/2605.10140v2. Full licence notice: \texttt{licenses/arxiv-2605.10140-CC-BY-4.0.txt}.\par\smallskip

\noindent\textit{Editorial scope.} Counted unit: the article's theorem thm:odd-coeff (source0.tex lines 1777--1799). The article's complete original proof of it is delivered with it (source0.tex lines 1801--1947, one \texttt{proof} environment of 147 lines). No mathematics is written, repaired or re-derived: the statement and the proof are the article's own lines, in the article's own notation, and no retained line is substituted or re-flowed.  Retained uncounted as the source-local context the proof needs: lines 1541--1560.  Nothing was omitted from any retained span, so the package carries no omission record.  No retained line contains a citation, so the wrapper carries no bibliography and no bibliography entry.  No retained line contains a figure environment, an image-inclusion command or drawing code, so no image is delivered and the package declares no asset dependency.  Editorial formatting only, in addition: an AMS \texttt{amsart} preamble that restates the article's own declarations and macros used by the retained lines, namely the environments \texttt{lemma}, \texttt{theorem} on separate counters, together with the standard packages those lines need.  The retained environments print in this document as Lemma 1, Theorem 1 in order of appearance, whereas the article prints the same items with the numbering of its own full document.  The complete native source of the article as distributed in the arXiv e-print for this exact version is bundled unchanged as \texttt{sources/200280/source0.tex}.
\begin{lemma}[Odd Hall averaging inequality]\label{lem:odd-Hall}
Let $\theta:\mathbb R\to\mathbb R$ be nondecreasing and satisfy
\[
        \theta(x+\pi)=\theta(x)+\pi .
\]
For $0\leq \tau\leq \pi/2$, define
\[
        J(\tau)
        =
        \frac1{2\pi}\int_0^{2\pi}
        \sin^2\frac{\theta(x+\tau)-\theta(x-\tau)}2\,\dd x .
\]
If $M$ is nonnegative and nonincreasing on $[0,\pi/2]$, then
\[
        \int_0^{\pi/2} M(\tau)J(\tau)\,\dd\tau
        \leq
        \frac2\pi
        \int_0^{\pi/2} M(\tau)\tau\,\dd\tau .
\]
\end{lemma}

\begin{theorem}[Odd boundary coefficient estimate]\label{thm:odd-coeff}
Let
\[
        f=h+\overline g
\]
be a sense-preserving harmonic diffeomorphism of $\D$ onto $\D$, with
\[
        h(z)=\sum_{n=0}^\infty a_nz^n,
        \qquad
        g(z)=\sum_{n=1}^\infty b_nz^n .
\]
Assume that $f$ extends to an orientation-preserving homeomorphism of
$\overline\D$ and that its boundary values are odd:
\[
        f(-\zeta)=-f(\zeta),
        \qquad |\zeta|=1 .
\]
Then
\[
        |a_1|^2+|b_1|^2\geq \frac{8}{\pi^2}.
\]
The constant is sharp.
\end{theorem}

\begin{proof}
Write the boundary values as
\[
        F(e^{\ii t})=f(e^{\ii t})=e^{\ii\theta(t)},
\]
where $\theta$ is a nondecreasing lift satisfying
\[
        \theta(t+2\pi)=\theta(t)+2\pi .
\]
Since the boundary map is odd, the lift may be chosen so that
\[
        \theta(t+\pi)=\theta(t)+\pi .
\]

Let
\[
        c_n=\frac1{2\pi}\int_0^{2\pi}
        F(e^{\ii t})e^{-\ii nt}\,\dd t
\]
be the Fourier coefficients of $F$. Since $F$ is odd, all even Fourier
coefficients vanish. Moreover, from the Poisson representation of
$f=h+\overline g$,
\[
        c_1=a_1,
        \qquad
        c_{-1}=\overline{b_1}.
\]
For $n\geq 1$, set
\[
        S_n=|c_n|^2+|c_{-n}|^2 .
\]

Hall's autocorrelation identity gives
\[
        C(t):=\frac1{2\pi}\int_0^{2\pi}
        \cos\bigl(\theta(s+t)-\theta(s-t)\bigr)\,\dd s
        =
        \sum_{n=1}^{\infty} S_n\cos(2nt).
\]
Since only odd Fourier modes occur,
\[
        C(t)=\sum_{k=0}^{\infty}
        S_{2k+1}\cos\bigl(2(2k+1)t\bigr).
\]
Therefore, by orthogonality on $[0,\pi/2]$,
\[
        S_1=
        \frac4\pi\int_0^{\pi/2} C(t)\cos(2t)\,\dd t .
\]
The odd-mode expansion also gives
\[
        C(\pi/2-t)=-C(t).
\]
Hence
\[
        S_1=
        \frac8\pi\int_0^{\pi/4} C(t)\cos(2t)\,\dd t .
\]

Define
\[
        J(t)=
        \frac1{2\pi}\int_0^{2\pi}
        \sin^2\frac{\theta(s+t)-\theta(s-t)}2\,\dd s .
\]
Then
\[
        C(t)=1-2J(t).
\]
Consequently,
\[
        S_1
        =
        \frac8\pi\int_0^{\pi/4}
        \bigl(1-2J(t)\bigr)\cos(2t)\,\dd t .
\]

Apply Lemma~\ref{lem:odd-Hall} with
\[
        M(t)=
        \begin{cases}
        \cos(2t), & 0\leq t\leq \pi/4,\\
        0, & \pi/4<t\leq \pi/2.
        \end{cases}
\]
This $M$ is nonnegative and nonincreasing on $[0,\pi/2]$. Hence
\[
        \int_0^{\pi/4} J(t)\cos(2t)\,\dd t
        \leq
        \frac2\pi\int_0^{\pi/4} t\cos(2t)\,\dd t .
\]
Therefore
\[
\begin{aligned}
        S_1
        &\geq
        \frac8\pi\int_0^{\pi/4}\cos(2t)\,\dd t
        -
        \frac{16}{\pi}\cdot\frac2\pi
        \int_0^{\pi/4}t\cos(2t)\,\dd t  \\
        &=
        \frac8{\pi^2}.
\end{aligned}
\]
Because
\[
        S_1=|c_1|^2+|c_{-1}|^2
        =|a_1|^2+|b_1|^2,
\]
the desired estimate follows.

It remains to show sharpness. Consider the four-point boundary collapse
\[
        F_\square(e^{\ii t})=\ii^k,
        \qquad
        \frac{k\pi}{2}\leq t<\frac{(k+1)\pi}{2},
        \qquad k=0,1,2,3.
\]
This map is not itself a boundary homeomorphism, but it is the
$L^1(\partial\D)$-limit of odd orientation-preserving circle
homeomorphisms. By the Radó--Kneser--Choquet theorem, the Poisson
extensions of these approximating homeomorphisms are sense-preserving
harmonic diffeomorphisms of $\D$ onto $\D$.

The Poisson extension of $F_\square$ has, up to rotation,
\[
        h'(z)=\frac{2(1-\ii)}{\pi}\frac1{1-z^4},
        \qquad
        g'(z)=-\frac{2(1+\ii)}{\pi}\frac{z^2}{1-z^4}.
\]
Thus
\[
        a_1=h'(0)=\frac{2(1-\ii)}{\pi},
        \qquad
        b_1=g'(0)=0.
\]
Consequently,
\[
        |a_1|^2+|b_1|^2
        =
        \left|\frac{2(1-\ii)}{\pi}\right|^2
        =
        \frac8{\pi^2}.
\]
The approximating odd boundary homeomorphisms therefore show that the
constant $8/\pi^2$ cannot be improved.
\end{proof}

\end{document}
